\chapter{Conclusions and Research Sequence}
\index{research sequence|textbf}

\label{ch:conclusions-and-outlook}

\chapterstatus{synthesis of the present evidence boundary and proposed work.
No completed production reduction or physical validation is claimed.}

The monograph proceeds from criteria for a good model, through the bulk-silicon
and doped-silicon representation programs, to the mathematical proof program
that records their obligations. The current project state does not yet supply a
material-specific numerical answer to the central reduction question. It does,
however, retain accepted synthetic numerical evidence for periodic reduction,
spin-space compatibility, and one-dimensional impurity extraction.

\section{Present conclusions}

Several conclusions are supported at the present evidence level.

First, spectral fitting and represented-operator fitting are different
requirements. Eigenvalues do not identify localized orbitals, gauge, site and
orbital correspondence, or the matrix elements through which later impurity
operators act. A reduced model intended only for spectral interpolation may be
adequate for that use while remaining unsuitable as a common coordinate system
for perturbations.

Second, pristine and doped matrix subtraction is not defined by equal array
shape. The retained state spaces, basis ordering, geometry, units, spin
convention, energy zero, and gauge must agree or be related by an explicit
identification map. A difference formed before these prerequisites are met has
no established impurity-operator meaning.

Third, projection, disentanglement, gauge transformation, Wannier localization,
and real-space truncation are separate operations. Their assumptions and losses
must remain attributable. In particular, localization does not itself establish
a low-rank approximation or authorize an unstated hopping cutoff.

Fourth, parent-model, numerical, and reduction errors answer different
questions. The PBE parent's relationship to experiment is not the same as a
finite-cutoff error, and neither is the same as the error introduced by an
$sp^3s^{\ast}$ model class. Combining them without a compatible mathematical
contract would hide rather than explain the quality of the reduction.

Fifth, the controlled periodic benchmarks establish a bounded verification
chain from one-dimensional plane-wave and finite-difference parents to
two-dimensional scalar, composite-band, topology, hopping, and localization
diagnostics. The two-dimensional optimizer study is a completed negative result
for its frozen synthetic protocol: its mesh-convergence and repeated-basin
criteria are unsupported. Because only one trial-projection family was
executed, the result establishes neither projection-family robustness, global
optimization, nor general Wannier convergence.

Sixth, the accepted spin-space and one-dimensional defect benchmarks verify
representation compatibility, folding, known-map and blind alignment,
independent extraction routes, finite-rank spectral checks, model-class
separation, finite-volume diagnostics, and independently refined continuum
comparisons for their declared synthetic systems. They do not establish a
two-dimensional defect operator, a silicon impurity model, or material
transferability.

Finally, the repository has an accepted represented-operator software
foundation, a coordinated proposed proof program, one frozen operator identity
checked by a pinned Lean backend, retained tutorial and bootstrap execution
observations, and human-accepted synthetic periodic and impurity-extraction
benchmarks. Those artifacts support their bounded proof, software, workflow,
and numerical-verification claims. They do not supply the converged production
parent, material Wannier operator, reduced bulk model, first-principles impurity
operator, or continuum result needed to answer the silicon questions.

\section{Outputs and perspective}

\subsection{Scholarly outputs}
\index{scholarly output}


The monograph remains the broad account of definitions, methods, provenance,
results, and limitations. Article manuscripts extract narrower questions.
P01, for example, can focus on bulk model-class expressiveness and the
compatibility of spectral and operator criteria without claiming completion of
the impurity program. Later articles may address alignment, impurity-operator
hierarchies, or continuum crossover only when their evidence exists.

Presentations may use the same retained results at a different explanatory
level, but they do not change evidentiary status. A conference figure derived
from tutorial data remains tutorial evidence; a proposed workflow diagram
remains proposed work; a production result retains its calculation identity and
limitations.

\subsection{Final perspective}

The intended contribution is not a single fitted parameter table. It is a
controlled account of when representations can be identified, when reduced
operators preserve their declared parent information, and where approximation
paths cease to agree. Reaching that account requires the full chain from
physical specification through numerical convergence, representation,
alignment, reduction, and evidence.

At present, the project has established the framework, bounded software
foundation, and a synthetic verification chain, but it has not completed the
proof packages or material reduction sequence. The accepted spin-space and
one-dimensional defect contracts and the accepted periodic two-dimensional
parent are frozen prerequisites for the two-dimensional defect extension in
Appendix~\ref{app:two-dimensional-defect-extraction}. That controlled synthetic
task is active but has no calculated result; it should test only genuinely
two-dimensional defect structure rather than repeat the completed
prerequisites. Material transfer follows only after that separate controlled
stage and its own authorization.

The remaining work is therefore stated as work: complete and review the
applicable proofs, construct the accepted material parent, measure each loss,
retain negative outcomes, and draw conclusions only at the strength supported
by the resulting evidence.
